\section{Results}\label{sec:results}

\subsection{Data quality and the confound floor}\label{sec:res-qc}
The fixed preprocessing produced \nWindows{} windows from all 88 subjects (Figure~\ref{fig:cohort});
ICA removed on average \icaExcludedMean{} components per recording. The expected EEG slowing in
dementia was present before any model was trained: the median theta/alpha power ratio was
\thetaAlphaAD{} in AD, \thetaAlphaFTD{} in FTD and \thetaAlphaCN{} in controls, with the posterior alpha
of controls replaced by diffuse delta and theta in AD (Figure~\ref{fig:topo}).

\begin{figure}[t]
\centering
\includegraphics[width=0.9\linewidth]{fig02_cohort}
\caption{Cohort: age, sex and recording length by group.}
\label{fig:cohort}
\end{figure}

\begin{figure}[t]
\centering
\includegraphics[width=0.9\linewidth]{fig03_band_topomaps}
\caption{Group-mean relative band power per electrode (rows: groups; columns: bands; colour scales
shared within each band).}
\label{fig:topo}
\end{figure}

A logistic regression on age and sex alone reached a macro-F1 of \demoFOne{}, against \chanceFOne{}
for a stratified random classifier. Demographics therefore carry some group information, consistent
with the sex imbalance between groups (Table~\ref{tab:demographics}). We use \demoFOne{}, not 1/3, as
the floor that EEG-based models must clear.

\subsection{Representation matters; graph construction does not}\label{sec:res-grid}
Across the 21 cells on the selection folds (\nMainRuns{} runs; Table~\ref{tab:cells}), the best
configuration was \bestCell{}, with subject-level macro-F1 \bestFOne{}~$\pm$~\bestFOneSd{}
(subject-bootstrap 95\% CI \bestFOneCiLo{}--\bestFOneCiHi{}). Because this cell was chosen as the
maximum of 21, that figure is optimistically biased. Re-evaluated on the independent confirmation
folds, the same cell scored \confFOne{}~$\pm$~\confFOneSd{} (95\% CI \confFOneCiLo{}--\confFOneCiHi{};
balanced accuracy \confBalAcc{}, macro ROC-AUC \confAuc{}, Cohen's $\kappa$ \confKappa{}) and again
ranked first of 21. The ordering of all cells was largely preserved (Spearman $\rho=\confRankRho$
between the two fold sets; Figure~\ref{fig:confirm}). We therefore take \confFOne{} as the
selection-free estimate for the best configuration.

Representation had a large effect on both fold sets (Table~\ref{tab:anova}; selection:
$F(\anovaRepDfOne,\anovaRepDfTwo)=\anovaRepF$, Greenhouse--Geisser $p\anovaRepP$, partial
$\eta^2=\anovaRepEta$; confirmation: $F=\canovaRepF$, $p\canovaRepP$, $\eta^2=\canovaRepEta$).
Averaged over graph types, the wavelet (P3, \repPThree{}; confirmation \crepPThree{}), spectral
(P1, \repPOne{}; \crepPOne{}) and fused (P6, \repPSix{}; \crepPSix{}) representations formed a leading
group; raw windows (P5, \repPFive{}) were intermediate; and connectivity strength (P2, \repPTwo{}),
PSWE features (P7, \repPSeven{}) and time-domain statistics (P4, \repPFour{}) trailed
(Figure~\ref{fig:rep}). Graph construction had a small effect on the selection folds
($F(\anovaEdgeDfOne,\anovaEdgeDfTwo)=\anovaEdgeF$, $p\anovaEdgeP$, partial $\eta^2=\anovaEdgeEta$;
marginal macro-F1 \edgespatial{} spatial, \edgefunctional{} functional, \edgehybrid{} hybrid;
Figure~\ref{fig:edge}) that did not replicate on the confirmation folds ($F=\canovaEdgeF$,
$p\canovaEdgeP$; marginal means \cedgespatial{}, \cedgefunctional{} and \cedgehybrid{}). The
representation~$\times$~graph interaction was not significant on either
(selection $F(\anovaInterDfOne,\anovaInterDfTwo)=\anovaInterF$, $p\anovaInterP$; confirmation
$F=\canovaInterF$, $p\canovaInterP$; Figure~\ref{fig:interaction}). A Friedman test across all cells
was significant ($p\friedmanP$), but the Nemenyi critical difference (\cdValue{} ranks) left the top
cells statistically indistinguishable from one another (Figure~\ref{fig:cd}).

Pairwise comparisons depend on how the dependence between cross-validation splits is handled. With
Holm-corrected Wilcoxon tests, \nPairsWilcoxonSig{} of the \nPairs{} representation pairs differed on
the selection folds and \cPairsWilcoxonSig{} on the confirmation folds; with the more conservative
Nadeau--Bengio corrected $t$-test, \nPairsNbSig{} and \cPairsNbSig{} did. For example, P3 exceeded P4
by \pThreeFourDiff{} macro-F1 on the selection folds (Cohen's $d=\pThreeFourD$; Wilcoxon--Holm
$p\pThreeFourWilP$; Nadeau--Bengio--Holm $p\pThreeFourNbP$). We regard the omnibus effect and the
ordering, which replicated, as robust, and individual pairwise differences as suggestive.

\begin{table}[t]
\centering\scriptsize
\caption{Subject-level performance of every cell (mean $\pm$ SD over the 25 selection splits,
seeds 0--4). CI: subject-bootstrap 95\% interval of macro-F1. BA: balanced accuracy; AUC: one-vs-rest
macro ROC-AUC. Confirm.: macro-F1 of the same cell on the 25 independent confirmation splits
(seeds 5--9).}
\label{tab:cells}
\setlength{\tabcolsep}{3.5pt}
\begin{tabular}{llccccccc}
\toprule
Representation & Graph & Macro-F1 & 95\% CI & BA & AUC & $\kappa$ & Brier & Confirm. \\
\midrule
\inputtab{generated/tab/cells.tex}
\bottomrule
\end{tabular}
\end{table}

\begin{table}[t]
\centering\small
\caption{Two-way repeated-measures ANOVA on subject-level macro-F1 (units: 25 test splits), on the
selection folds and, independently, on the confirmation folds. $\varepsilon$: Greenhouse--Geisser;
$p$: GG-corrected; $\eta^2_p$: partial eta squared.}
\label{tab:anova}
\setlength{\tabcolsep}{4pt}
\begin{tabular}{lccccccccc}
\toprule
 & & \multicolumn{4}{c}{Selection (seeds 0--4)} & \multicolumn{4}{c}{Confirmation (seeds 5--9)} \\
\cmidrule(lr){3-6}\cmidrule(lr){7-10}
Effect & df & $F$ & $\varepsilon$ & $p$ & $\eta^2_p$ & $F$ & $\varepsilon$ & $p$ & $\eta^2_p$ \\
\midrule
\inputtab{generated/tab/anova.tex}
\bottomrule
\end{tabular}
\end{table}

\begin{figure}[t]
\centering
\includegraphics[width=0.62\linewidth]{fig06_representation}\hfill
\includegraphics[width=0.33\linewidth]{fig07_edge}
\caption{Subject-level macro-F1 by representation (left, averaged over graph types) and by graph
construction (right, averaged over representations). Points are the 25 test splits; black markers,
mean $\pm$ 95\% CI; dotted line, 1/3.}
\label{fig:rep}\label{fig:edge}
\end{figure}

\begin{figure}[t]
\centering
\includegraphics[width=\linewidth]{fig08_interaction}
\caption{Mean macro-F1 of every representation~$\times$~graph cell (left) and the same values as an
interaction plot (right). Lines are close to parallel: the best representation does not depend on
how the graph is built.}
\label{fig:interaction}
\end{figure}

\begin{figure}[t]
\centering
\includegraphics[width=\linewidth]{fig09_critical_difference}
\caption{Critical-difference diagram over the 21 cells. Cells joined by a black bar are not separated
by the Nemenyi test.}
\label{fig:cd}
\end{figure}

\subsection{What the models get right and wrong}\label{sec:res-errors}
For \bestCell{}, sensitivity was \bestSensCN{} for controls, \bestSensAD{} for AD and only
\bestSensFTD{} for FTD (specificities \bestSpecCN{}, \bestSpecAD{} and \bestSpecFTD{}). Across
representations, controls were recognised best and FTD --- the smallest group --- worst; FTD
recordings were most often misclassified as AD, except with PSWE features, where they were more
often taken for controls (Figure~\ref{fig:confusion}). The ROC and precision--recall
curves show the same ordering (Figure~\ref{fig:roc}). The models were poorly calibrated (expected
calibration error \bestEce{}), so their probabilities should not be read as clinical risks.

\begin{figure}[t]
\centering
\includegraphics[width=\linewidth]{fig10_confusion_matrices}
\caption{Row-normalised confusion matrices summed over the 25 test splits, for the best graph type
of each representation.}
\label{fig:confusion}
\end{figure}

\begin{figure}[t]
\centering
\includegraphics[width=\linewidth]{fig11_roc_pr}
\caption{One-vs-rest ROC (top) and precision--recall (bottom) curves for the three best cells; mean
$\pm$ SD across the five repeats.}
\label{fig:roc}
\end{figure}

\subsection{Graph models do not beat their non-graph twins}\label{sec:res-baselines}
Table~\ref{tab:baselines} compares every representation's best GCN cell with classical classifiers
trained on the same features, test splits and weighting, under both training regimes. With training
subjects only, the best classical model reached \classicalTrBest{}; adding the validation subjects
raised it to \classicalTvBest{}, and refitting the GCN on the same subjects raised \bestCell{} to
\refitBestFOne{} (mean gain over all cells \refitMeanGain{}). The paired difference between each
representation's best GCN cell and its best non-graph twin ranged from \fairTrDeltaMin{} to
\fairTrDeltaMax{} in the train-only regime and from \fairTvDeltaMin{} to \fairTvDeltaMax{} in the
train + val regime, and none was significant under either Holm-corrected test (Table~\ref{tab:fair};
for P3, train only: \fairTrPThreeDelta{}, Wilcoxon--Holm $p\fairTrPThreeWilP$; train + val:
\fairTvPThreeDelta{}, $p\fairTvPThreeWilP$). The GCN on wavelet features was numerically the best
model in both regimes, but its advantage over non-graph classifiers is within the noise of this
cohort. The spectral, wavelet, fused and raw representations cleared the demographic floor of
\demoFOne{}; the connectivity (P2), time-domain (P4) and PSWE (P7) representations scored within
about 0.05 of it.

\begin{table}[t]
\centering\small
\caption{Every representation's best GCN cell against its non-graph twins (same test splits,
weighting and subject aggregation) under both training regimes, and the demographic floor.
Macro-F1, mean $\pm$ SD over the 25 selection splits. Train only: models fitted on the training
subjects (the GCN early-stopped on the validation subjects). Train + val: models fitted on training
and validation subjects (the GCN refitted for its early-stopped epoch count). GCN graph type: best
in each regime.}
\label{tab:baselines}
\begin{tabular}{llcc}
\toprule
Representation & Model & Train only & Train + val \\
\midrule
\inputtab{generated/tab/baselines.tex}
\bottomrule
\end{tabular}
\end{table}

\subsection{Ablations}\label{sec:res-ablations}
Table~\ref{tab:ablations} and Figure~\ref{fig:forest} report each ablation on the confirmation folds
as a paired difference from the same cell's confirmation-grid run. No ablation reached significance
under either Holm-corrected test (\cablNSigWil{} of 26 with Wilcoxon, \cablNSigNb{} with
Nadeau--Bengio), and every point estimate except one was negative or zero: no modification improved
on the reference configuration. Removing the graph altogether (identity adjacency) changed
macro-F1 by \cablPThreespatialAOnenographDelta{} [\cablPThreespatialAOnenographLo{},
\cablPThreespatialAOnenographHi{}] for the spatial cell and by \cablPThreehybridAOnenographDelta{}
for the hybrid cell, and a degree-preserving random graph by \cablPThreespatialATworandomgraphDelta{}
and \cablPThreehybridATworandomgraphDelta{}. Any benefit of message passing is therefore small, and
we cannot attribute it to a specific electrode topology rather than to smoothing over neighbours.
Two changes had uncorrected $t$-intervals that excluded zero for the spatial cell: GAT attention
(\cablPThreespatialASixgatDelta{} [\cablPThreespatialASixgatLo{}, \cablPThreespatialASixgatHi{}]) and a
fully connected graph (\cablPThreespatialAThreefullDelta{} [\cablPThreespatialAThreefullLo{},
\cablPThreespatialAThreefullHi{}]), the latter consistent with over-smoothing; neither interval
accounts for the overlap between cross-validation training sets. Sparsity ($k=3$ or 6), binarised
edges, depth and loss weighting changed macro-F1 by at most about 0.02. Regressing age and sex out of
every feature cost \cablPThreespatialAOneTworesidualizedDelta{} [\cablPThreespatialAOneTworesidualizedLo{},
\cablPThreespatialAOneTworesidualizedHi{}] for the spatial cell, leaving performance well above the
demographic floor. The selection-fold ablations (Table~\ref{tab:ablations-main}) gave the same
conclusions.

\begin{table}[t]
\centering\scriptsize
\caption{Ablations on the confirmation folds: paired $\Delta$ macro-F1 (ablated minus the same
cell's confirmation-grid run) with 95\% $t$-interval over the 25 splits, and Nadeau--Bengio $p$ with
Holm correction within each cell. Selection-fold values: Table~\ref{tab:ablations-main}.}
\label{tab:ablations}
\begin{tabular}{lcccc}
\toprule
 & \multicolumn{2}{c}{P3 $\times$ spatial} & \multicolumn{2}{c}{P3 $\times$ hybrid} \\
\cmidrule(lr){2-3}\cmidrule(lr){4-5}
Ablation & $\Delta$ [95\% CI] & $p_\mathrm{Holm}$ & $\Delta$ [95\% CI] & $p_\mathrm{Holm}$ \\
\midrule
\inputtab{generated/tab/ablations.tex}
\bottomrule
\end{tabular}
\end{table}

\begin{figure}[t]
\centering
\includegraphics[width=0.75\linewidth]{fig13_ablation_forest}
\caption{Ablations on the confirmation folds as paired differences from the two reference cells
(mean and 95\% CI over 25 splits). An asterisk would mark Nadeau--Bengio--Holm $p<0.05$; none qualifies.}
\label{fig:forest}
\end{figure}

\subsection{Leakage inflates results to the published range}\label{sec:res-leak}
Re-running the two best cells on the confirmation folds with window-level splitting --- identical
data, features, graph, model and hyperparameters --- raised subject-level macro-F1 from
\cleakPThreespatialCorrectSubj{} to \cleakPThreespatialLeakySubj{} (spatial) and from
\cleakPThreehybridCorrectSubj{} to \cleakPThreehybridLeakySubj{} (hybrid). Window-level macro-F1, the
metric most often reported, rose to \cleakPThreespatialLeakyWin{} and \cleakPThreehybridLeakyWin{}
(Figure~\ref{fig:leak}). The selection folds gave almost identical values (\leakPThreespatialCorrectSubj{}
to \leakPThreespatialLeakySubj{} for the spatial cell). The leaky protocol thus reproduces the
95--99\% figures reported on this benchmark with a model that, evaluated correctly, scores about 0.6.

\begin{figure}[t]
\centering
\includegraphics[width=0.8\linewidth]{fig12_leakage}
\caption{The same pipeline evaluated with subject-level (correct) and window-level (leaky)
splitting on the confirmation folds. Bars: mean $\pm$ SD over 25 splits; dotted line, 1/3.}
\label{fig:leak}
\end{figure}

\subsection{The best cell is above chance}\label{sec:res-perm}
On the first confirmation repeat, the observed macro-F1 of \bestCell{} was \cpermPThreespatialObs{},
against a label-permutation null with mean \cpermPThreespatialNullMean{} and maximum
\cpermPThreespatialNullMax{} over \cpermPThreespatialN{} permutations ($p=\cpermPThreespatialP$, the
smallest value attainable; Figure~\ref{fig:perm}). The original test on the selection folds gave the
same result (\permPThreespatialN{} permutations, observed \permPThreespatialObs{}, null maximum
\permPThreespatialNullMax{}, $p=\permPThreespatialP$).

\subsection{PSWE: a BBB-associated marker that tracks slowing}\label{sec:res-pswe}
Fixed-threshold PSWEs were most frequent in AD (median \psweRateAD{} events/min), intermediate in FTD
(\psweRateFTD{}) and least frequent in controls (\psweRateCN{}; Kruskal--Wallis $p\psweKwP$;
Figures~\ref{fig:pswe-example} and~\ref{fig:pswe}). In negative-binomial models adjusted for age and
sex, the PSWE rate ratio relative to controls was \psweAdjADRR{} (95\% CI \psweAdjADLo{}--\psweAdjADHi{},
$p\psweAdjADP$) in AD and \psweAdjFTDRR{} (\psweAdjFTDLo{}--\psweAdjFTDHi{}, $p\psweAdjFTDP$) in FTD.

The PSWE rate was, however, almost collinear with background slowing (Spearman $\rho=\psweSlowRho$
with relative delta+theta power; age- and sex-adjusted partial $\rho=\psweSlowPartial$). Once slowing
was added to the model, the AD rate ratio reversed to \psweAdjSlowADRR{} (\psweAdjSlowADLo{}--\psweAdjSlowADHi{},
$p\psweAdjSlowADP$), and FTD fell to \psweAdjSlowFTDRR{} ($p\psweAdjSlowFTDP$). The PSWE rate was not
significantly associated with MMSE within patients ($\rho=\psweMmseRho$, $p\psweMmseP$). In the
classifier, PSWE features alone (P7) performed near the demographic floor (\repPSeven{}), and adding
them to the best representation, or using PSWE co-occurrence as the graph, did not help
(Section~\ref{sec:res-ablations}).

\paragraph{Background-relative PSWEs (exploratory).} Defining PSWEs relative to each channel's own
background reversed the group pattern. With the primary definition (a drop of at least \SI{2}{\hertz}
below the channel median), AD showed \emph{fewer} transient slowing events than controls (rate ratio
\reldropTwoAdjADRR{}, \reldropTwoAdjADLo{}--\reldropTwoAdjADHi{}, $p\reldropTwoAdjADP$), FTD a
non-significant reduction (\reldropTwoAdjFTDRR{}), and the event rate was only weakly related to
slowing ($\rho=\reldropTwoSlowRho$) and not to MMSE ($\rho=\reldropTwoMmseRho$). After additional
adjustment for background MPF the AD effect was no longer significant (\reldropTwoAdjBgADRR{},
$p\reldropTwoAdjBgADP$). The MAD-scaled variants gave the same direction more strongly and survived
that adjustment (for 2~MAD: AD \relmadTwoAdjBgADRR{}, $p\relmadTwoAdjBgADP$; FTD \relmadTwoAdjBgFTDRR{},
$p\relmadTwoAdjBgFTDP$). Every variant pointed the same way: relative to their own background,
patients with dementia had fewer, not more, transient slowing events.

\begin{figure}[t]
\centering
\includegraphics[width=\linewidth]{fig04_pswe_example}
\caption{Fixed-threshold PSWE detection in the AD recording with the highest PSWE rate among the
example shard (first 120\,s, channel with most events). Top: EEG; bottom: median power frequency
with the \SI{6}{\hertz} threshold; shaded, detected PSWEs. Because this subject's background is
already slow, most of the recording is flagged --- the pattern behind the collinearity with slowing.}
\label{fig:pswe-example}
\end{figure}

\begin{figure}[t]
\centering
\includegraphics[width=\linewidth]{fig05_pswe_burden}
\caption{Fixed-threshold PSWE rate per subject by group (left), median per-channel rate as scalp maps
(middle), and rate against MMSE in patients (right; analysis only).}
\label{fig:pswe}
\end{figure}
